\subsection{An unsupervised Bayesian neural network model for phylodynamics}

\begin{figure}
    \centering
    \includegraphics[width=0.7\linewidth]{figures/BELLA.pdf}
    \caption{\rv{Schematic representation of the proposed \textit{Bayesian evolutionary layered learning architectures} (BELLA) phylodynamic framework. The model takes as input time-stamped sequence data (or a fixed, time-calibrated phylogenetic tree), together with predictor variables including, but not limited to, time \(\mathbf{t}\), time series covariates \(\mathbf{x}\), and traits $\mathbf{T} = (T^{(j)})_{j=1}^N$. BELLA flexibly maps these predictors to phylodynamic parameters---such as speciation, extinction, transmission, or migration rates---allowing for complex, nonlinear relationships that generalize classical generalized linear models (GLMs) and skyline-type approaches. For example, in macroevolution \(T^{(j)}\) may represent a phenotypic trait, with the output corresponding to the diversification rate associated with species possessing that trait for each time point. All parameters, including neural network weights, are estimated jointly within a Markov chain Monte Carlo framework. The resulting posterior distributions enable both inference of population dynamics and interpretation of predictor effects using explainable artificial intelligence tools.}}
    \label{fig:bella}
\end{figure}

We implemented an unsupervised Bayesian neural network to model the variation of phylodynamic parameters as a function of one or multiple predictors, which we termed \emph{Bayesian Evolutionary Layered Learning Architectures} (BELLA, Fig.~\ref{fig:bella}). The BELLA model operates within the phylogenetic inference framework BEAST 2 \cite{bouckaert2019beast} that can jointly infer tree topology, branch lengths, the parameters of evolutionary substitution models, as well as phylodynamic parameters directly from alignments of molecular and morphological data.
The model applies similarly to epidemiology, where the focal phylodynamic parameters may include transmission, recovery, and migration rates \citep{Vaughan2025BDMM}, and to macroevolution, where speciation, extinction, and fossilization rates describe species diversification and sampling within a fossilized birth--death (FBD) process \citep{Heath2014FBD}. The predictors used as input to the neural network include time itself or other time series---such as those describing environmental variation---as well as lineage-specific features, for example phenotypic traits in macroevolutionary processes or population mobility data in epidemiological contexts.

In the implemented framework, we employ a phylodynamic birth–death model \cite{stadler2010sampling} to describe the branching of the tree, in which the phylodynamic parameters are the birth, death, migration, and sampling rates, determined by predictors through time and across lineages.
While our framework is described here in the context of the birth--death model, the same concepts can also be applied to other phylodynamic approaches, such as the coalescent \cite{kingman1982coalescent}.

As in standard phylodynamic analyses, the BELLA model relies on unsupervised Bayesian inference to sample the parameters from their posterior probability distribution using a Markov chain Monte Carlo (MCMC) algorithm.
Yet, in contrast to most implementations, the BELLA model does not sample the parameters of the phylodynamic process directly; instead, it samples the weights of the neural network, from which such phylodynamic parameters are obtained.
Owing to its representational flexibility, the neural network serves as a proxy for the potentially complex mapping between input predictors and phylodynamic parameters, enabling the model to capture nonlinear, non-monotonic responses and interactions among predictors.
